% ============================================================
%  BANK SOAL MATEMATIKA SMA — KURIKULUM MERDEKA
%  BAGIAN 2: FASE F  —  VERSI SOAL (REVISI)  |  No. 66–100
%  - Tanpa banner; opsi A-E tersusun ke bawah
%  - Sinkron dengan Versi Lengkap (soal & nomor identik)
%  Kelas dokumen: exam  |  Kompilasi: pdfLaTeX (Overleaf)
% ============================================================
\documentclass[11pt,a4paper]{exam}

\usepackage[utf8]{inputenc}
\usepackage[T1]{fontenc}
\usepackage[bahasa]{babel}
\usepackage{amsmath,amssymb}
\usepackage{geometry}
\usepackage{xcolor}
\usepackage{tikz}
\usepackage{pgfplots}
\pgfplotsset{compat=1.18}
\usepackage{enumitem}
\usepackage{array,booktabs}

\geometry{a4paper,margin=2.3cm}

% (Versi Soal: TANPA \printanswers)

\definecolor{biru}{HTML}{1F4E79}
\definecolor{hijau}{HTML}{2E7D32}
\definecolor{oranye}{HTML}{C75B12}
\definecolor{ungu}{HTML}{6A1B9A}

\pagestyle{headandfoot}
\runningheadrule
\firstpageheadrule
\header{\textbf{\textcolor{biru}{BANK SOAL MATEMATIKA}}}{}{\textbf{FASE F}}
\footer{}{Halaman \thepage}{}

\renewcommand{\choicelabel}{\Alph{choice}.}

\renewcommand{\choiceshook}{%
  \setlength{\leftmargin}{2em}%
  \setlength{\itemsep}{5pt}%
  \setlength{\topsep}{5pt}%
  \setlength{\parsep}{0pt}}

\newcommand{\Fund}{{\footnotesize\textsf{\textcolor{hijau}{[Fundamental]}}}\ }
\newcommand{\Inter}{{\footnotesize\textsf{\textcolor{biru}{[Intermediate]}}}\ }
\newcommand{\Adv}{{\footnotesize\textsf{\textcolor{oranye}{[Advanced]}}}\ }
\newcommand{\Exp}{{\footnotesize\textsf{\textcolor{ungu}{[Expert]}}}\ }

\newcommand{\topik}[1]{%
  \par\addvspace{14pt}%
  \noindent{\large\bfseries\color{biru}#1}\par\addvspace{6pt}}

\begin{document}

\begin{center}
{\color{biru}\rule{\linewidth}{1.5pt}}\\[4pt]
{\LARGE\bfseries\color{biru} BANK SOAL MATEMATIKA SMA}\\[2pt]
{\large Versi Soal — Bagian 2: Fase F}\\[2pt]
{\small 35 Soal Pilihan Ganda (No. 66–100)}\\[4pt]
{\color{biru}\rule{\linewidth}{1.5pt}}
\end{center}
\vspace{4pt}

\noindent\textit{Petunjuk:} Pilih satu jawaban yang paling tepat. Label tingkat
kesulitan tertera pada awal tiap soal.

\setcounter{question}{65}

\topik{9. Komposisi Fungsi (7 Soal)}

\begin{questions}

\question \Fund Diketahui $f(x) = 2x+1$ dan $g(x) = x^{2}$. Hasil $(f\circ g)(x)$ adalah \dots
\begin{choices}
\choice $(2x+1)^{2}$ \choice $4x^{2}+1$ \CorrectChoice $2x^{2}+1$
\choice $2x^{2}+2x+1$ \choice $x^{2}+1$
\end{choices}

\question \Fund Diketahui $f(x) = x+3$ dan $g(x) = 2x$. Hasil $(f\circ g)(x)$ adalah \dots
\begin{choices}
\CorrectChoice $2x+3$ \choice $2x+6$ \choice $x+6$ \choice $2x$ \choice $6x$
\end{choices}

\question \Inter Diketahui $f(x) = x+3$ dan $g(x) = 2x$. Nilai $(g\circ f)(2)$ adalah \dots
\begin{choices}
\choice $7$ \choice $8$ \CorrectChoice $10$ \choice $5$ \choice $16$
\end{choices}

\question \Inter Diketahui $f(x) = 2x-1$ dan $g(x) = x+4$. Hasil $(f\circ g)(x)$ adalah \dots
\begin{choices}
\choice $2x+3$ \choice $2x+8$ \choice $2x-1$ \CorrectChoice $2x+7$ \choice $x+7$
\end{choices}

\question \Inter Diketahui $(f\circ g)(x) = 4x+5$ dan $g(x) = 2x+1$. Maka $f(x)$ adalah \dots
\begin{choices}
\CorrectChoice $2x+3$ \choice $2x+1$ \choice $4x+3$ \choice $2x+5$ \choice $x+3$
\end{choices}

\question \Adv Diketahui $f(x) = x^{2}+1$ dan $g(x) = x-2$. Nilai $(f\circ g)(3)$ adalah \dots
\begin{choices}
\CorrectChoice $2$ \choice $10$ \choice $8$ \choice $5$ \choice $17$
\end{choices}

\question \Adv Diketahui $f(x) = 2x+1$ dan $(g\circ f)(x) = 6x+7$. Maka $g(x)$ adalah \dots
\begin{choices}
\choice $3x+1$ \choice $6x+4$ \choice $3x+7$ \choice $2x+4$ \CorrectChoice $3x+4$
\end{choices}

\topik{10. Fungsi Invers (7 Soal)}

\question \Fund Invers dari fungsi $f(x) = 2x+1$ adalah \dots
\begin{choices}
\choice $\dfrac{x+1}{2}$ \CorrectChoice $\dfrac{x-1}{2}$ \choice $2x-1$
\choice $\dfrac{1-x}{2}$ \choice $\dfrac{x}{2}-1$
\end{choices}

\question \Fund Invers dari fungsi $f(x) = x-3$ adalah \dots
\begin{choices}
\choice $x-3$ \choice $3-x$ \choice $\dfrac{1}{x-3}$ \CorrectChoice $x+3$ \choice $3x$
\end{choices}

\question \Inter Invers dari fungsi $f(x) = \dfrac{x+2}{3}$ adalah \dots
\begin{choices}
\choice $3x+2$ \choice $\dfrac{x-2}{3}$ \CorrectChoice $3x-2$
\choice $3(x-2)$ \choice $\dfrac{x}{3}-2$
\end{choices}

\question \Inter Diketahui $f(x) = 2x-4$. Nilai $f^{-1}(6)$ adalah \dots
\begin{choices}
\choice $8$ \CorrectChoice $5$ \choice $6$ \choice $10$ \choice $2$
\end{choices}

\question \Inter Diketahui $f(x) = \dfrac{3x-2}{x+1}$, $x \neq -1$. Maka $f^{-1}(x)$ adalah \dots
\begin{choices}
\CorrectChoice $\dfrac{x+2}{3-x}$ \choice $\dfrac{x-2}{x+3}$ \choice $\dfrac{3x-2}{x+1}$
\choice $\dfrac{x+1}{3x-2}$ \choice $\dfrac{2-x}{x-3}$
\end{choices}

\question \Adv Diketahui $f(x) = x^{2}+1$ untuk $x \ge 0$. Maka $f^{-1}(x)$ adalah \dots
\begin{choices}
\choice $\sqrt{x+1}$ \choice $x^{2}-1$ \CorrectChoice $\sqrt{x-1}$
\choice $\sqrt{x}-1$ \choice $(x-1)^{2}$
\end{choices}

\question \Adv Diketahui $f(x) = 2x+3$ dan $g(x) = x-1$. Maka $(f\circ g)^{-1}(x)$ adalah \dots
\begin{choices}
\choice $\dfrac{x+1}{2}$ \CorrectChoice $\dfrac{x-1}{2}$ \choice $\dfrac{x-3}{2}$
\choice $2x-1$ \choice $x-1$
\end{choices}

\topik{11. Lingkaran (10 Soal)}

\question \Fund Persamaan lingkaran dengan pusat $O(0,0)$ dan jari-jari 5 adalah \dots
\begin{choices}
\choice $x^{2}+y^{2}=5$ \choice $x+y=5$ \CorrectChoice $x^{2}+y^{2}=25$
\choice $x^{2}+y^{2}=10$ \choice $(x-5)^{2}+y^{2}=25$
\end{choices}

\question \Fund Pusat lingkaran $(x-3)^{2}+(y+2)^{2}=16$ adalah \dots
\begin{choices}
\choice $(-3,2)$ \CorrectChoice $(3,-2)$ \choice $(3,2)$ \choice $(-3,-2)$ \choice $(2,-3)$
\end{choices}

\question \Fund Jari-jari lingkaran $(x-1)^{2}+(y-2)^{2}=49$ adalah \dots
\begin{choices}
\choice $49$ \choice $14$ \choice $\sqrt{7}$ \CorrectChoice $7$ \choice $24{,}5$
\end{choices}

\question \Inter Persamaan lingkaran berpusat $(2,-1)$ dan berjari-jari 3 adalah \dots
\begin{choices}
\CorrectChoice $(x-2)^{2}+(y+1)^{2}=9$
\choice $(x+2)^{2}+(y-1)^{2}=9$
\choice $(x-2)^{2}+(y+1)^{2}=3$
\choice $(x-2)^{2}+(y-1)^{2}=9$
\choice $(x+2)^{2}+(y+1)^{2}=9$
\end{choices}

\question \Inter Pusat lingkaran $x^{2}+y^{2}-4x+2y-4=0$ adalah \dots
\begin{choices}
\choice $(-2,1)$ \CorrectChoice $(2,-1)$ \choice $(4,-2)$ \choice $(2,1)$ \choice $(-4,2)$
\end{choices}

\question \Inter Jari-jari lingkaran $x^{2}+y^{2}-4x+2y-4=0$ adalah \dots
\begin{choices}
\choice $9$ \choice $\sqrt{8}$ \CorrectChoice $3$ \choice $4$ \choice $\sqrt{6}$
\end{choices}

\question \Inter Perhatikan juring lingkaran berikut dengan jari-jari 6 cm dan
sudut pusat $60^\circ$.

\begin{center}
\begin{tikzpicture}[scale=0.9]
\draw[thick,biru] (0,0) circle (1.8);
\fill[biru!15] (0,0)--(1.8,0) arc (0:60:1.8)--cycle;
\draw[biru,thick] (0,0)--(1.8,0);
\draw[biru,thick] (0,0)--(60:1.8);
\node at (1.1,0.45){\small $60^\circ$};
\node[below] at (0.9,0){\small $r=6$};
\end{tikzpicture}
\end{center}

Panjang busur juring tersebut adalah \dots
\begin{choices}
\choice $6\pi$ cm \choice $\pi$ cm \choice $12\pi$ cm \CorrectChoice $2\pi$ cm \choice $3\pi$ cm
\end{choices}

\question \Adv Luas juring lingkaran berjari-jari 10 cm dengan sudut pusat
$72^\circ$ adalah \dots
\begin{choices}
\choice $10\pi$ cm$^2$ \CorrectChoice $20\pi$ cm$^2$ \choice $72\pi$ cm$^2$
\choice $50\pi$ cm$^2$ \choice $25\pi$ cm$^2$
\end{choices}

\question \Adv Posisi titik $(5,1)$ terhadap lingkaran $x^{2}+y^{2}-4x+2y-4=0$ adalah \dots
\begin{choices}
\choice pada lingkaran
\choice di dalam lingkaran
\CorrectChoice di luar lingkaran
\choice di pusat lingkaran
\choice tidak dapat ditentukan
\end{choices}

\question \Exp Persamaan garis singgung lingkaran $x^{2}+y^{2}=25$ di titik $(3,4)$ adalah \dots
\begin{choices}
\CorrectChoice $3x+4y=25$ \choice $4x+3y=25$ \choice $3x+4y=5$
\choice $3x-4y=25$ \choice $4x+3y=5$
\end{choices}

\topik{12. Statistika Bivariat (11 Soal)}

\question \Fund Data dua variabel (bivariat) paling tepat disajikan dengan \dots
\begin{choices}
\choice histogram
\CorrectChoice diagram pencar (scatter plot)
\choice diagram lingkaran
\choice box plot
\choice diagram batang
\end{choices}

\question \Fund Jika titik-titik pada diagram pencar cenderung naik dari kiri bawah
ke kanan atas, korelasinya bersifat \dots
\begin{choices}
\CorrectChoice positif \choice negatif \choice nol \choice tidak ada \choice konstan
\end{choices}

\question \Inter Nilai koefisien korelasi $r$ selalu berada pada rentang \dots
\begin{choices}
\choice $0 \le r \le 1$ \choice $r \ge 0$ \CorrectChoice $-1 \le r \le 1$
\choice $-1 \le r \le 0$ \choice $0 \le r \le 100$
\end{choices}

\question \Inter Nilai $r$ yang mendekati $-1$ menunjukkan \dots
\begin{choices}
\choice korelasi positif kuat
\choice tidak ada korelasi
\choice korelasi positif lemah
\CorrectChoice korelasi negatif kuat
\choice data konstan
\end{choices}

\question \Inter Perhatikan diagram pencar berikut.

\begin{center}
\begin{tikzpicture}
\begin{axis}[width=7.5cm,height=4.8cm,xlabel=$x$,ylabel=$y$,
  xmin=0,xmax=10,ymin=0,ymax=10,grid=both,grid style={black!12}]
\addplot[only marks,mark=*,biru,mark size=1.6pt]coordinates{
(1,9)(2,8)(3,8)(4,6)(5,6)(6,5)(7,3)(8,3)(9,2)};
\end{axis}
\end{tikzpicture}
\end{center}

Pola hubungan yang ditunjukkan adalah \dots
\begin{choices}
\choice korelasi positif \CorrectChoice korelasi negatif \choice tidak ada korelasi
\choice korelasi positif sempurna \choice data konstan
\end{choices}

\question \Inter Tabel berikut mencatat lama belajar dan nilai ujian beberapa siswa.

\begin{center}\small
\begin{tabular}{|c|c|c|c|c|}
\hline
Jam belajar ($x$) & 1 & 2 & 3 & 4\\\hline
Nilai ujian ($y$) & 60 & 70 & 80 & 90\\\hline
\end{tabular}
\end{center}

Variabel ``jam belajar'' berperan sebagai \dots
\begin{choices}
\CorrectChoice variabel bebas \choice variabel terikat \choice konstanta
\choice pencilan \choice median
\end{choices}

\question \Adv Garis tren data (jam belajar $x$, nilai $y$) adalah $\hat{y} = 5x + 50$.
Prediksi nilai siswa yang belajar 6 jam adalah \dots
\begin{choices}
\choice $75$ \CorrectChoice $80$ \choice $85$ \choice $90$ \choice $50$
\end{choices}

\question \Adv Perhatikan diagram pencar berikut.

\begin{center}
\begin{tikzpicture}
\begin{axis}[width=7.5cm,height=4.8cm,xlabel=$x$,ylabel=$y$,
  xmin=0,xmax=6,ymin=0,ymax=10,grid=both,grid style={black!12}]
\addplot[only marks,mark=*,biru,mark size=1.8pt]coordinates{
(1,2)(2,3)(3,5)(4,6)(5,7)};
\addplot[only marks,mark=*,oranye,mark size=2.2pt]coordinates{(3,9)};
\end{axis}
\end{tikzpicture}
\end{center}

Titik yang merupakan pencilan (\emph{outlier}) adalah \dots
\begin{choices}
\choice $(1,2)$ \CorrectChoice $(3,9)$ \choice $(5,7)$ \choice $(2,3)$ \choice $(4,6)$
\end{choices}

\question \Adv Sebuah penelitian menemukan korelasi positif kuat antara penjualan
es krim dan jumlah kasus tenggelam. Kesimpulan yang paling tepat adalah \dots
\begin{choices}
\choice penjualan es krim menyebabkan orang tenggelam
\choice kasus tenggelam menyebabkan orang membeli es krim
\CorrectChoice keduanya dipengaruhi faktor ketiga (cuaca panas); korelasi bukan sebab-akibat
\choice data penelitian pasti salah
\choice tidak ada hubungan apa pun
\end{choices}

\question \Adv Perhatikan tabel data berpasangan berikut.

\begin{center}\small
\begin{tabular}{|c|c|c|c|c|}
\hline
$x$ & 1 & 2 & 3 & 4\\\hline
$y$ & 2 & 4 & 6 & 8\\\hline
\end{tabular}
\end{center}

Jenis hubungan kedua variabel tersebut adalah \dots
\begin{choices}
\choice korelasi negatif
\choice tidak ada korelasi
\CorrectChoice korelasi positif sempurna ($r=1$)
\choice korelasi positif lemah
\choice $r=0$
\end{choices}

\question \Exp Perhatikan diagram pencar hubungan umur mobil (tahun) dan harga jual
(juta rupiah) berikut.

\begin{center}
\begin{tikzpicture}
\begin{axis}[width=7.5cm,height=4.8cm,xlabel=Umur (tahun),ylabel=Harga (juta),
  xmin=0,xmax=7,ymin=80,ymax=200,grid=both,grid style={black!12}]
\addplot[only marks,mark=*,oranye,mark size=1.8pt]coordinates{
(1,180)(2,165)(3,150)(4,138)(5,122)(6,108)};
\end{axis}
\end{tikzpicture}
\end{center}

Perkiraan nilai koefisien korelasi yang paling tepat adalah \dots
\begin{choices}
\CorrectChoice $\approx -0{,}9$ \choice $\approx 0{,}9$ \choice $\approx 0{,}1$
\choice $\approx -0{,}3$ \choice $0$
\end{choices}

\end{questions}

\vspace{6pt}
\begin{center}
{\color{biru}\rule{\linewidth}{1pt}}\\[3pt]
{\small\textit{Akhir Bagian 2 (Fase F) — Versi Soal. Kunci jawaban \& pembahasan ada pada Versi Lengkap.}}
\end{center}

\end{document}